\documentclass[10pt,a4paper]{amsart}
\usepackage[margin=19mm]{geometry}
\usepackage{amsmath,amssymb,mathtools}
\usepackage[scr=rsfs,cal=euler]{mathalfa}
\usepackage{xfrac}
\usepackage[unicode,hidelinks,bookmarks=false]{hyperref}
\newcommand{\ie}{i.e.\ }
\newcommand{\cf}{cf.\ }
\newcommand{\resp}{respectively\ }
\newcommand{\Subsection}{\subsection}
\providecommand{\nobreakdash}{\nobreak\hbox{-}}
\setlength{\emergencystretch}{2em}
\multlinegap0pt
\usepackage[shortlabels]{enumitem}
\usepackage{amssymb}
\usepackage{mathtools}
\usepackage{tikz}
\newtheorem{lemma}{Lemma}
\newcommand{\Aset}{\mathsf A}
\newcommand{\Bset}{\mathsf B}
\newcommand{\Eset}{\mathsf E}
\newcommand{\toq}{\to_q}
\title{The Quantum Homomorphism Orders of Graphs Are Universal}
\author{Yangjing Long}
\date{Source: arXiv:2605.19543v2 (CC BY 4.0)}
\begin{document}
\maketitle

\noindent\emph{Source and licence.} The Quantum Homomorphism Orders of Graphs Are Universal. Version arXiv:2605.19543v2 (CC BY 4.0), distributed by arXiv under the Creative Commons Attribution 4.0 International licence, http://creativecommons.org/licenses/by/4.0/, as recorded in the arXiv licence metadata for this exact version and shown on the abstract page https://arxiv.org/abs/2605.19543v2. Full licence notice: \texttt{licenses/arxiv-2605.19543-CC-BY-4.0.txt}.\par\smallskip

\noindent\textit{Editorial scope.} Counted unit: the article's lemma lem:color-localization (source0.tex lines 405--412). The article's complete original proof of it is delivered with it (source0.tex lines 414--530, one \texttt{proof} environment of 117 lines). No mathematics is written, repaired or re-derived: the statement and the proof are the article's own lines, in the article's own notation, and no retained line is substituted or re-flowed.  Retained uncounted as the source-local context the proof needs: lines 376--383.  Nothing was omitted from any retained span, so the package carries no omission record.  No retained line contains a citation, so the wrapper carries no bibliography and no bibliography entry.  No retained line contains a figure environment, an image-inclusion command or drawing code, so no image is delivered and the package declares no asset dependency.  Editorial formatting only, in addition: an AMS \texttt{amsart} preamble that restates the article's own declarations and macros used by the retained lines, namely the environments \texttt{lemma} on separate counters, together with the standard packages those lines need.  The retained environments print in this document as Lemma 1 in order of appearance, whereas the article prints the same items with the numbering of its own full document.  The complete native source of the article as distributed in the arXiv e-print for this exact version is bundled unchanged as \texttt{sources/200922/source0.tex}.
\begin{lemma}\label{lem:structural-color}
Let $Y=H*J$ for an irreflexive directed graph $H$.  Then:
\begin{enumerate}[label=\textup{(\roman*)}]
\item If $y\notin \Aset(Y)$, then the induced graph $Y[N_Y(y)]$ is bipartite.
\item If $y\notin \Bset(Y)$, then $Y[\Aset(Y)\cap N_Y(y)]$ contains no edge.
\item If $y\notin \Eset(Y)$, then $Y[\Bset(Y)\cap N_Y(y)]$ contains no edge.
\end{enumerate}
\end{lemma}

\begin{lemma}[Color localization]\label{lem:color-localization}
Let $X$ be a graph containing vertices carrying the color gadgets above, let $Y=H*J$, and let $F:X\toq Y$ be a quantum homomorphism.  Then:
\begin{enumerate}[label=\textup{(\roman*)}]
\item If $x$ is an $A$-root in $X$ and $y\notin \Aset(Y)$, then $F_{x,y}=0$.
\item If $x$ is a $B$-root in $X$ and $y\notin \Bset(Y)$, then $F_{x,y}=0$.
\item If $x$ is an $E$-root in $X$ and $y\notin \Eset(Y)$, then $F_{x,y}=0$.
\end{enumerate}
\end{lemma}

\begin{proof}
We prove the three claims in order.

\smallskip
\noindent\emph{Proof of (i).}
Let $x$ be an $A$-root and suppose $y\notin \Aset(Y)$.  Put $P=F_{x,y}$.  The vertex $x$ is adjacent to the vertices $c_0,\ldots,c_4$ of a private $5$-cycle.  By Lemma~\ref{lem:structural-color}, $Y[N_Y(y)]$ is bipartite; fix a bipartition
\[
   N_Y(y)=N_0\sqcup N_1.
\]
For $i\in\mathbb Z_5$, define
\[
   A_i=\sum_{z\in N_0}F_{c_i,z},\qquad
   B_i=\sum_{z\in N_1}F_{c_i,z},\qquad
   C_i=A_i+B_i,\qquad S_i=A_i-B_i.
\]
Because the projections $F_{c_i,z}$, $z\in V(Y)$, form a projective measurement, the operators $A_i$ and $B_i$ are orthogonal projections, $C_i$ is a projection, and
\[
   S_i^2=C_i. \tag{1}
\]
Since $xc_i$ is an edge of $X$, the edge-zero relation gives
\[
   F_{c_i,z}P=0\qquad (z\notin N_Y(y)).
\]
Equivalently, since the projections are self-adjoint, also $PF_{c_i,z}=0$ for $z\notin N_Y(y)$.  Together with completeness for the measurement at $c_i$, this gives
\[
   C_iP=P\qquad\text{for every }i. \tag{2}
\]
Because $N_0$ and $N_1$ are independent sets in $Y[N_Y(y)]$, and because $c_ic_{i+1}$ is an edge of $X$, the edge-zero relation gives
\[
   A_iA_{i+1}=0,
   \qquad
   B_iB_{i+1}=0
   \qquad (i\in\mathbb Z_5). \tag{3}
\]
Using (2) and (3),
\[
\begin{aligned}
   S_iS_{i+1}P
   &=(A_i-B_i)(A_{i+1}-B_{i+1})P\\
   &=-(A_iB_{i+1}+B_iA_{i+1})P\\
   &=-C_iC_{i+1}P.
\end{aligned}
\]
Here
\[
   C_iC_{i+1}P=C_i(C_{i+1}P)=C_iP=P,
\]
where we use (2) and associativity of operator multiplication only; no commutativity of $C_i$ and $C_{i+1}$ is being assumed.  Hence
\[
   S_iS_{i+1}P=-P. \tag{4}
\]
The same argument with $i$ and $i+1$ interchanged gives
\[
   S_{i+1}S_iP=-P. \tag{5}
\]
Let $v\in\operatorname{Ran}(P)$.  Then (1) and (2) imply
\[
   \|S_iv\|^2=\langle S_i^2v,v\rangle=\langle C_iv,v\rangle=\|v\|^2.
\]
Equations (4) and (5) imply
\[
   (S_i+S_{i+1})^2v
   =(S_i^2+S_iS_{i+1}+S_{i+1}S_i+S_{i+1}^2)v=0.
\]
The operators $S_i$ and $S_{i+1}$ are self-adjoint, since they are differences of orthogonal projections.  Hence $S_i+S_{i+1}$ is self-adjoint and $(S_i+S_{i+1})^2$ is positive semidefinite.  Taking the inner product with $v$ gives
\[
   0=\langle (S_i+S_{i+1})^2v,v\rangle
    =\|(S_i+S_{i+1})v\|^2.
\]
Thus
\[
   (S_i+S_{i+1})v=0,
\]
so $S_{i+1}v=-S_iv$ for every $i\in\mathbb Z_5$.  Going around the odd cycle gives
\[
   S_0v=-S_4v=-S_0v,
\]
and hence $S_0v=0$.  On the other hand,
\[
   \|S_0v\|^2=\langle S_0^2v,v\rangle=\langle C_0v,v\rangle=\|v\|^2,
\]
where the last equality uses $C_0P=P$ and $v\in\operatorname{Ran}(P)$.  Thus $v=0$.  Therefore $\operatorname{Ran}(P)=0$, and $P=0$.

\smallskip
\noindent\emph{Proof of (ii).}
Let $x$ be a $B$-root and suppose $y\notin \Bset(Y)$.  Put $P=F_{x,y}$.  By construction, $x$ has two adjacent $A$-neighbors $\alpha_1,\alpha_2$.

For $i=1,2$, define
\[
   C_i=\sum_{z\in \Aset(Y)\cap N_Y(y)}F_{\alpha_i,z}.
\]
We claim that $C_iP=P$.  Indeed, by completeness at $\alpha_i$,
\[
   P=\sum_{z\in V(Y)}F_{\alpha_i,z}P.
\]
If $z\notin\Aset(Y)$, then $F_{\alpha_i,z}=0$ by part (i).  If $z\in\Aset(Y)$ but $z\notin N_Y(y)$, then $F_{\alpha_i,z}P=0$ because $x\alpha_i$ is an edge of $X$ while $yz$ is not an edge of $Y$.  Hence only terms with $z\in\Aset(Y)\cap N_Y(y)$ remain, and $C_iP=P$.

Since $y\notin\Bset(Y)$, Lemma~\ref{lem:structural-color} says that $\Aset(Y)\cap N_Y(y)$ contains no edge.  As $\alpha_1\alpha_2$ is an edge of $X$, the edge-zero condition gives
\[
   C_1C_2=0.
\]
Consequently
\[
   P=C_1P=C_1C_2P=0.
\]

\smallskip
\noindent\emph{Proof of (iii).}
Let $x$ be an $E$-root and suppose $y\notin\Eset(Y)$.  Put $P=F_{x,y}$.  By construction, $x$ has two adjacent $B$-neighbors $\beta_1,\beta_2$.  For $i=1,2$, define
\[
   D_i=\sum_{z\in \Bset(Y)\cap N_Y(y)}F_{\beta_i,z}.
\]
The same argument as in (ii), now using part (ii) in place of part (i), gives $D_iP=P$ for $i=1,2$.  Since $y\notin\Eset(Y)$, Lemma~\ref{lem:structural-color} says that $\Bset(Y)\cap N_Y(y)$ contains no edge.  Since $\beta_1\beta_2$ is an edge of $X$, the edge-zero condition gives $D_1D_2=0$.  Hence
\[
   P=D_1P=D_1D_2P=0.
\]
\end{proof}

\end{document}
